\section{Материальное обеспечение группы}

\begin{table}[h!]
	\centering
%	\resizebox{0.77\textwidth}{!}{%
		\begin{tabular}{|>{\centering\arraybackslash}m{0.02\linewidth}|>{\centering\arraybackslash}m{0.31\linewidth}|>{\centering\arraybackslash}m{0.08\linewidth}|>{\centering\arraybackslash}m{0.29\linewidth}|>{\centering\arraybackslash}m{0.08\linewidth}|}
			\hline
			\multirow{2}{*}{\textnumero}	&	\multicolumn{2}{|c|}{Общественное}	&	\multicolumn{2}{|c|}{Личное}	\\
			
			\cline{2-5} & Наименование	&	Кол-во	&	Наименование	&	Кол-во\\
			\hline
			1	&	Верёвка статическая 50~м, $\varnothing=9$~мм	&	1	&	Каска	&	1	\\
			\hline
			2	&	Карабины	&	4	&	Ледоруб	&	1	\\
			\hline
			3	&	Жумар	&	1	&	Солнцезащитные очки	&	1	\\
			\hline
			4	&	Страховочная система	&	2	& Треккинговые палки &	2 \\
			\hline
			5	&	Блокировка	&	2	& 	&	 	\\
			\hline
			6	&	Спусковое устройство	&	1	&	 	&	 	\\
			\hline
			7	&	Станционные петли	&	2	&	 	&	 \\
			\hline
			8	&  Кордалет	& 2		&		&	\\
			\hline
			9	& Спутниковый телефон	&	1	&	 	&	 \\
			\hline
			10 & Рация* & 1 & &\\ 
			\hline

		\end{tabular}%
%	}
\\
* Для связи с группой Алексея Остапива, который шёл параллельным маршрутом
\end{table}

Более подробное распределение веса общественного снаряжение приведен в таблице ниже.


\begin{tabular}{{|>{\centering\arraybackslash}m{0.15\linewidth}
			|>{\centering\arraybackslash}m{0.4\linewidth}
			|>{\centering\arraybackslash}m{0.15\linewidth}
			|>{\centering\arraybackslash}m{0.1\linewidth}|}}
	\hline
	Группа & Состав & Количество, шт & Вес, г \\
	\hline
	Кухня & Горелка & 2 & 250х2 \\
	\hline
	& Кан & 2 & 600х2 \\
	\hline
	& Шторки для канов от ветра & 2 & 550 \\
	\hline
	& Огнеупорная ткань + перчатки & 1 & 200 \\
	\hline
	& Поварешка & 1 & 100 \\
	\hline
	& Разделочная доска & 1 & 200 \\
	\hline
	& Губка для посуды & 1 & 30 \\
	\hline
	Горючее & Газовые баллоны 450г & 9 & 5976 \\
	\hline
\end{tabular}


\begin{tabular}{{|>{\centering\arraybackslash}m{0.15\linewidth}
			|>{\centering\arraybackslash}m{0.4\linewidth}
			|>{\centering\arraybackslash}m{0.15\linewidth}
			|>{\centering\arraybackslash}m{0.1\linewidth}|}}
	\hline
	Группа & Состав & Количество, шт & Вес, г \\
	\hline
	Бивак & Палатка четверка & 1 & 4000 \\
	\hline
	& Палатка двойка & 1 & 2500 \\
	\hline
	& Тент & 1 & 500 \\
	\hline
	& Пленка для рюкзаков & 1 & 300 \\
	\hline
	& Мешки для мусора, рулон & 1 & 100 \\
	\hline
	& Скотч обычный, малярный, красный & 3 & 190 \\
	\hline
	& Безмен & 1 & 60 \\
	\hline
	Веревки & Веревка 40м & 1 & 2900 \\
	\hline
	Специальное снаряжение & Карабины(4шт), жумар, спусковуха & 1 & 530 \\
	\hline
	& Страховочная система & 2 & 800х2 \\
	\hline
	& Блокировка & 2 & 300х2 \\
	\hline
	& Станционные петли(120см) & 2 & 300 \\
	\hline
	Ремнабор & См отчет реммастера & 1 & 2500 \\
	\hline
	Аптечка & См отчет медика & 1 & 3000 \\
	\hline
	Аппаратура & Спутниковый телефон & 1 & 300 \\
	\hline
	& Рация & 1 & 300 \\
	\hline
	& Фотоаппарат & 1 & 1300 \\
	\hline
	Прочее & Флаг ГС & 1 & 90 \\
	\hline
\end{tabular}
\begin{itemize}
	\item Общий вес общественного снаряжения : 29,826кг.
	\item	Специальное снаряжение – веревка, репшнур, петли, обвязки и карабины были взяты с целью обеспечить минимальный страховочный комплект в случае аварийной ситуации. В течение похода оно не применялись.
	\item В течении всего маршрута участники пользовались треккинговыми палками, а также надевали каски на камнеопасных склонах и на переходах рек вброд. 
	\item Тент использовался по назначению 3 августа на мн1  и 13 августа на месте обеда. Также 5 августа тент был использван для  установки палатки под сильным градом на мн3. 
	\item Расход газа составил 38г/чел/день.
\end{itemize}


 



\clearpage